\section{Conclusion}
\label{sec:conclusion}

This work implemented and evaluated two least-squares procedures using a
reduced QR factorization computed by the Modified Gram--Schmidt method.
The numerical routines for inner products, norms, matrix operations, QR
factorization, and backward substitution were written explicitly in Python
rather than delegated to high-level linear algebra solvers. The results
demonstrate how the QR formulation converts each overdetermined problem into
an upper-triangular system while allowing the computation to be assessed
through independent numerical indicators.

For the first problem, the relative QR reconstruction error was
approximately $1.07\times10^{-16}$, and the orthogonality error was
approximately $4.52\times10^{-16}$. The optimality indicator,
$\|A^T(A\mathbf{x}^*-\mathbf{b})\|_2$, was approximately
$6.10\times10^{-14}$. Together with the full-column-rank diagnostic, these
values are consistent with a numerically accurate and unique least-squares
solution for the matrix considered. The residual norm remained nonzero, as
expected for an inconsistent overdetermined system; it measures the distance
from $\mathbf b$ to the column space of $A$ rather than a failure of the
algorithm.

For the second problem, the estimated coefficients were
$a\approx1.97\times10^{-15}$, $b=0.4$, and $c=0.6$. Thus, although a
polynomial of degree at most two was prescribed, the fitted polynomial is
effectively linear at the reported precision: $p(x)\approx0.4x+0.6$. The
mean squared error was approximately $0.5333$, while the QR reconstruction,
orthogonality, and least-squares optimality indicators were approximately
$7.30\times10^{-17}$, $3.95\times10^{-15}$, and $2.31\times10^{-14}$,
respectively. This illustrates the distinction between numerical accuracy
and fit quality: small algorithmic errors do not imply that the model
reproduces every observation exactly.

The full-column-rank condition ensures uniqueness in both problems. For the
quadratic fit, this follows theoretically because the six distinct
abscissas make the three columns of the design matrix linearly independent.
Numerically, the diagonals of the computed triangular factors also exceeded
the adopted rank tolerance. Nevertheless, Modified Gram--Schmidt without
column pivoting is not a universally reliable rank-revealing method for
severely ill-conditioned matrices; a pivoted QR factorization or singular
value decomposition may be preferable when stronger rank diagnostics are
required \cite{higham2002accuracy}.

Finally, the observations in this assignment are generated from the
enrollment digits according to the prescribed rules. Consequently, the
second experiment illustrates polynomial least-squares fitting and
numerical validation, not an empirical inference about predator--prey
dynamics. Overall, the two experiments demonstrate the practical
connection between least-squares theory, explicit algorithmic
implementation, and numerical verification.
